\documentclass[10pt,a4paper]{amsart}
\usepackage{amsmath,amssymb,amsthm,graphicx}
\usepackage{verbatim}
\newenvironment{reference}{\comment}{\endcomment}
%\newenvironment{reference}{}{}
\newenvironment{slogan}{\comment}{\endcomment}
\newenvironment{history}{\comment}{\endcomment}

% For commutative diagrams we use Xy-pic
\usepackage[all]{xy}

% We use 2cell for 2-commutative diagrams.
\xyoption{2cell}
% Theorem environments.
%
\theoremstyle{plain}
\newtheorem{theorem}[subsection]{Theorem}
\newtheorem{proposition}[subsection]{Proposition}
\newtheorem{lemma}[subsection]{Lemma}

\theoremstyle{definition}
\newtheorem{definition}[subsection]{Definition}
\newtheorem{example}[subsection]{Example}
\newtheorem{exercise}[subsection]{Exercise}
\newtheorem{situation}[subsection]{Situation}

\theoremstyle{remark}
\newtheorem{remark}[subsection]{Remark}
\newtheorem{remarks}[subsection]{Remarks}

\numberwithin{equation}{subsection}

% Macros
%
\def\lim{\mathop{\mathrm{lim}}\nolimits}
\def\colim{\mathop{\mathrm{colim}}\nolimits}
\def\Spec{\mathop{\mathrm{Spec}}}
\def\Hom{\mathop{\mathrm{Hom}}\nolimits}
\def\Ext{\mathop{\mathrm{Ext}}\nolimits}
\def\SheafHom{\mathop{\mathcal{H}\!\mathit{om}}\nolimits}
\def\SheafExt{\mathop{\mathcal{E}\!\mathit{xt}}\nolimits}
\def\Sch{\mathit{Sch}}
\def\Mor{\mathop{\mathrm{Mor}}\nolimits}
\def\Ob{\mathop{\mathrm{Ob}}\nolimits}
\def\Sh{\mathop{\mathit{Sh}}\nolimits}
\def\NL{\mathop{N\!L}\nolimits}
\def\CH{\mathop{\mathrm{CH}}\nolimits}
\def\proetale{{pro\text{-}\acute{e}tale}}
\def\etale{{\acute{e}tale}}
\def\QCoh{\mathit{QCoh}}
\def\Ker{\mathop{\mathrm{Ker}}}
\def\Im{\mathop{\mathrm{Im}}}
\def\Coker{\mathop{\mathrm{Coker}}}
\def\Coim{\mathop{\mathrm{Coim}}}

% Boxtimes
%
\DeclareMathSymbol{\boxtimes}{\mathbin}{AMSa}{"02}

%
% Macros for moduli stacks/spaces
%
\def\QCohstack{\mathcal{QC}\!\mathit{oh}}
\def\Cohstack{\mathcal{C}\!\mathit{oh}}
\def\Spacesstack{\mathcal{S}\!\mathit{paces}}
\def\Quotfunctor{\mathrm{Quot}}
\def\Hilbfunctor{\mathrm{Hilb}}
\def\Curvesstack{\mathcal{C}\!\mathit{urves}}
\def\Polarizedstack{\mathcal{P}\!\mathit{olarized}}
\def\Complexesstack{\mathcal{C}\!\mathit{omplexes}}
% \Pic is the operator that assigns to X its picard group, usage \Pic(X)
% \Picardstack_{X/B} denotes the Picard stack of X over B
% \Picardfunctor_{X/B} denotes the Picard functor of X over B
\def\Pic{\mathop{\mathrm{Pic}}\nolimits}
\def\Picardstack{\mathcal{P}\!\mathit{ic}}
\def\Picardfunctor{\mathrm{Pic}}
\def\Deformationcategory{\mathcal{D}\!\mathit{ef}}
\usepackage{cleveref}
\crefname{originallemma}{original lemma}{original lemmas}
\Crefname{originallemma}{Original lemma}{Original lemmas}
\makeatletter
\newlabel{lemma-one-ideal-ci}{{0B03}{source}{}{sourcepointer-1}{}}
\newlabel{lemma-one-ideal-ci@cref}{{[originallemma][0][]0B03}{[1][0][]source}{}{}{}}
\makeatother
\crefname{originalreference}{original reference}{original references}
\Crefname{originalreference}{Original reference}{Original references}
\makeatletter
\newlabel{algebra-proposition-dimension}{{00KQ}{source}{}{sourcepointer-2}{}}
\newlabel{algebra-proposition-dimension@cref}{{[originalreference][0][]00KQ}{[1][0][]source}{}{}{}}
\makeatother
\crefname{originalreference}{original reference}{original references}
\Crefname{originalreference}{Original reference}{Original references}
\makeatletter
\newlabel{algebra-lemma-dimension-prime-polynomial-ring}{{00P0}{source}{}{sourcepointer-3}{}}
\newlabel{algebra-lemma-dimension-prime-polynomial-ring@cref}{{[originalreference][0][]00P0}{[1][0][]source}{}{}{}}
\makeatother
\crefname{originalreference}{original reference}{original references}
\Crefname{originalreference}{Original reference}{Original references}
\makeatletter
\newlabel{divisors-lemma-Noetherian-scheme-regular-ideal}{{063I}{source}{}{sourcepointer-4}{}}
\newlabel{divisors-lemma-Noetherian-scheme-regular-ideal@cref}{{[originalreference][0][]063I}{[1][0][]source}{}{}{}}
\makeatother
\crefname{originallemma}{original lemma}{original lemmas}
\Crefname{originallemma}{Original lemma}{Original lemmas}
\makeatletter
\newlabel{lemma-pullback-by-regular-immersion}{{0AZN}{source}{}{sourcepointer-5}{}}
\newlabel{lemma-pullback-by-regular-immersion@cref}{{[originallemma][0][]0AZN}{[1][0][]source}{}{}{}}
\makeatother
\crefname{originalreference}{original reference}{original references}
\Crefname{originalreference}{Original reference}{Original references}
\makeatletter
\newlabel{algebra-lemma-dimension-at-a-point-finite-type-field}{{00P1}{source}{}{sourcepointer-6}{}}
\newlabel{algebra-lemma-dimension-at-a-point-finite-type-field@cref}{{[originalreference][0][]00P1}{[1][0][]source}{}{}{}}
\makeatother
\title{LCI intersection multiplicity as leading coefficient of a Hilbert polynomial}
\author{The Stacks Project Authors}
\date{stacks-project:intersection.tex:local-snapshot-20261012; GNU FDL 1.2 or later}
\begin{document}
\maketitle
\noindent\textit{Editorial note.} All delivered mathematical text is the author's own, copied verbatim from the pinned Stacks Project chapter file intersection.tex; no statement or proof step was rewritten, reordered or removed.\par\medskip
\noindent\textit{Editorial note.} Span recomputation by label pairing (LANE\_BRIEF 10.2): in this chapter the author's proof environment follows the statement environment as a standalone block; the spec covers the complete environment pair (statement span, proof span) as recorded in the pool candidate.\par\medskip
\noindent\textit{Editorial note.} Occupation scan: combined claim\_key + statement\_anchor + generic-line scan over all registration files (active and withdrawn) found 0 hits for this chapter (is-occupied.json, 2026-10-14); consistent with the duty manager's six-check dispatch.\par\medskip
\ifdefined\hypertarget\hypertarget{sourcepointer-1}{}\fi
\noindent\textit{Source reference.} Original lemma 0B03: sources/stacks-chow/tags.txt, source lines 13846--13846. Cross-chapter standard prerequisite cited by the author inside the counted span; it is used as a pointer and its own proof is not part of this item.\par\medskip
\ifdefined\hypertarget\hypertarget{sourcepointer-2}{}\fi
\noindent\textit{Source reference.} Original reference 00KQ: sources/stacks-chow/tags.txt, source lines 760--760. Cross-chapter standard prerequisite cited by the author inside the counted span; it is used as a pointer and its own proof is not part of this item.\par\medskip
\ifdefined\hypertarget\hypertarget{sourcepointer-3}{}\fi
\noindent\textit{Source reference.} Original reference 00P0: sources/stacks-chow/tags.txt, source lines 914--914. Cross-chapter standard prerequisite cited by the author inside the counted span; it is used as a pointer and its own proof is not part of this item.\par\medskip
\ifdefined\hypertarget\hypertarget{sourcepointer-4}{}\fi
\noindent\textit{Source reference.} Original reference 063I: sources/stacks-chow/tags.txt, source lines 7832--7832. Cross-chapter standard prerequisite cited by the author inside the counted span; it is used as a pointer and its own proof is not part of this item.\par\medskip
\ifdefined\hypertarget\hypertarget{sourcepointer-5}{}\fi
\noindent\textit{Source reference.} Original lemma 0AZN: sources/stacks-chow/tags.txt, source lines 13831--13831. Cross-chapter standard prerequisite cited by the author inside the counted span; it is used as a pointer and its own proof is not part of this item.\par\medskip
\ifdefined\hypertarget\hypertarget{sourcepointer-6}{}\fi
\noindent\textit{Source reference.} Original reference 00P1: sources/stacks-chow/tags.txt, source lines 915--915. Cross-chapter standard prerequisite cited by the author inside the counted span; it is used as a pointer and its own proof is not part of this item.\par\medskip
\begin{lemma}
\label{lemma-multiplicity-with-lci}
Let $X$ be a nonsingular variety. Let $W,V \subset X$ be
closed subvarieties which intersect properly. Let $Z$ be an irreducible
component of $V \cap W$ with generic point $\xi$.
Suppose the ideal of $V$ in $\mathcal{O}_{X, \xi}$ is cut out by
a regular sequence $f_1, \ldots, f_c \in \mathcal{O}_{X, \xi}$.
Then $e(X, V\cdot W, Z)$ is equal to $c!$ times the leading coefficient in
the Hilbert polynomial
$$
t \mapsto \text{length}_{\mathcal{O}_{X, \xi}}
\mathcal{O}_{W, \xi}/(f_1, \ldots, f_c)^t,\quad t \gg 0.
$$
In particular, this coefficient is $> 0$.
\end{lemma}

\begin{proof}
The equality
$$
e(X, V\cdot W, Z) = e_{(f_1, \ldots, f_c)}(\mathcal{O}_{W, \xi}, c)
$$
follows from the more general Lemma \ref{lemma-one-ideal-ci}.
To see that $e_{(f_1, \ldots, f_c)}(\mathcal{O}_{W, \xi}, c)$ is
$> 0$ or equivalently that $e_{(f_1, \ldots, f_c)}(\mathcal{O}_{W, \xi}, c)$
is the leading coefficient of the Hilbert polynomial
it suffices to show that the
dimension of $\mathcal{O}_{W, \xi}$ is $c$, because the degree of the
Hilbert polynomial is equal to the dimension by
Algebra, Proposition \ref{algebra-proposition-dimension}.
Say $\dim(V) = r$, $\dim(W) = s$, and $\dim(X) = n$. Then
$\dim(Z) = r + s - n$ as the intersection is proper. Thus
the transcendence degree of $\kappa(\xi)$ over $\mathbf{C}$ is
$r + s - n$, see Algebra, Lemma
\ref{algebra-lemma-dimension-prime-polynomial-ring}.
We have $r + c = n$ because $V$ is cut out by a regular sequence
in a neighbourhood of $\xi$, see
Divisors, Lemma \ref{divisors-lemma-Noetherian-scheme-regular-ideal}
and then Lemma \ref{lemma-pullback-by-regular-immersion}
applies (for example). Thus
$$
\dim(\mathcal{O}_{W, \xi}) = s - (r + s - n) = s - ((n - c) + s - n) = c
$$
the first equality by Algebra, Lemma
\ref{algebra-lemma-dimension-at-a-point-finite-type-field}.
\end{proof}

\end{document}
